\section{Related Work}
\label{sec:relat}

\textbf{Novel View Synthesis.}
Novel View Synthesis generates unseen scene perspectives from limited input images. Neural Radiance Fields (NeRF)~\cite{mildenhall2021nerf} pioneered photorealistic synthesis through volumetric rendering with implicit neural representations. Subsequent work advanced three directions: geometry enhancement~\cite{wang2021neus}, rendering quality~\cite{barron2021mipnerf,barron2022mipnerf360}, and accelerated training and rendering via hybrid representations such as hash grids and voxels~\cite{muller2022instantngp,chen2022tensorf}.
3D Gaussian Splatting (3DGS)~\cite{kerbl20233dgs} revolutionized NVS with anisotropic 3D Gaussians and tile-based rasterization, achieving state-of-the-art quality and speed. Follow-ups further improved geometric fidelity through surface regularization~\cite{guedon2024sugar, lu2024scaffold}, anti-aliasing~\cite{yu2024mipsplatting}, and texture and gradient optimization~\cite{yu2024gof, chen2024pgsr}. Addressing 3DGS's geometric inconsistencies, 2D Gaussian Splatting (2DGS)~\cite{huang20242dgs} projects disks onto explicit surfaces with local smoothing for view-consistent reconstruction. Our work adopts 2D Gaussian primitives to improve geometric accuracy.

\noindent\textbf{Inverse Rendering of Glossy Objects.}
Inverse rendering estimates scene geometry, materials, and lighting from images but suffers from material-light ambiguities and complex light interactions on glossy surfaces. NeRF-based approaches employ neural radiance fields with ray marching but incur computational bottlenecks from dense sampling. Ref-NeRF~\cite{verbin2022refnerf} simplifies view-dependent effects via directional encoding, yet omits material-light decomposition. TensoIR~\cite{jin2023tensoir} uses tensor factorization for geometry and ray marching for indirect illumination, while ENVIDR~\cite{liang2023envidr} employs surface-based modeling of glossy reflections via a decomposed neural renderer. NeRO~\cite{liu2023nero} reconstructs geometry using split-sum approximation without masks, utilizing integrated directional encoding for illumination, then recovers BRDF and lighting via Monte Carlo sampling. Despite flexibility, these implicit methods incur significant overhead from dense ray sampling.

Recent 3DGS-based methods significantly advance inverse rendering. Pioneering works~\cite{jiang2024gaussianshader, liang2024gsir} apply simplified per-primitive rendering to Gaussians but face geometric and material fidelity limitations. Subsequent methods demonstrate divergent strengths: 3DGS-DR~\cite{ye20243dgsdr} optimizes high-frequency details via deferred shading; Ref-Gaussian~\cite{yao2024refgaussian} uses split-sum approximation for glossy surfaces, inspiring our geometry reconstruction. Parallel developments integrate specialized priors such as microfacet segmentation~\cite{lai2025glossygs}, foundation model supervision~\cite{tong2025gs2dgs}, and diffusion models~\cite{du2024gsid} to reduce geometry ambiguities, motivating our geometry-aware initialization. R3DG~\cite{gao2024r3dg} models the full rendering equation, extended by IRGS~\cite{gu2025irgs} with differentiable 2D Gaussian ray tracing~\cite{Moenne-3DGSRT} for precise visibility and indirect illumination, a technique adopted in our pipeline. Finally, spherical mipmap encoding and directional factorization~\cite{zhang2025refgs, kouros2025rgsdr} enhance specular effects, guiding our specular compensation.